%----------------------------------------------------------
% 제7장. 사례 연구: 한국 기업의 AI 거버넌스 실천
%----------------------------------------------------------

\chapter{Case 연구: 한국 기업의 AI 거버넌스 실천}
\label{ch:case_studies}

이론과 프레임워크는 실제 적용을 통해 비로소 의미를 갖는다. 
본 장에서는 한국 기업과 공공기관\cite{enholm2022artificial}\cite{kuziemski2020ai}의 AI 거버넌스 구축 사례를 분석한다. 
각 사례는 도입 배경, 거버넌스 체계 설계, 직면한 도전과 해결 방안, 
성능와 교훈을 포함한다.

본 장의 사례는 공개된 정보, 업계 Report\cite{ransbotham2017reshaping}, 저자의 프로젝트 경험을 바탕으로 
재구성한 것이다. 일부 세부 사항은 기밀 보호를 위해 일반화하거나 
유사 사례를 종합하여 제시하였다.


%----------------------------------------------------------
\section{금융 사례: KBank의 AI 신용평가 거버넌스}
\label{sec:7.1}

\subsection{도입 배경}
\label{sec:7.1.1}

KBank은 2022년 AI based 신용평가 시스템을 도입\cite{dixon2020machine}하면서 
포괄적인 AI 거버넌스 체계를 구축하였다. 
금융 산업의 특성상 규제 준수와 고객 신뢰\cite{bartlett2022consumer}가 핵심 고려사항이었다.

\paragraph{도입 동기}

\begin{itemize}
 \item 기존 통계 based 신용평가 모델의 한계 (thin-file 고객 Response 어려움)
 \item 핀테크 경쟁사 대비 심사 속도 개선\cite{philippon2020fintech} 필요
 \item 대안 Data(Alternative 데이터) 활용을 통한 금융 포용성\cite{cao2022ai} 확대
 \item 금융위원회 AI 가이드라인\cite{fsc2024ai}(2021) Response
\end{itemize}

\paragraph{초기 도전}

\begin{table}[htbp]
\centering
\caption{KBank AI 신용평가 도입 시 주요 도전}
\label{tab:kbank_challenges}
\begin{tabularx}{\textwidth}{p{3cm}X}
\toprule
\textbf{영역} & \textbf{도전 과제} \\
\midrule
규제 준수 & 
신용정보법 상 금지 변수 정의 불명확; 대안 데이터의 적법성 결정 기준 부재 \\
\midrule
공정성 & 
역사적 대출 데이터에 내재된 편향\cite{buolamwini2018gender}; 연령.지역별 승인율 격차 우려 \\
\midrule
설명가능성\cite{arrieta2020explainable} & 
딥러닝 모델의 블랙박스 특성; 고객 거절 사유 설명 의무 \\
\midrule
내부 역량 & 
AI 윤리 전문 인력 부재; 기존 MRM(Model 리스크 관리) 체계와 AI 거버넌스 통합 필요 \\
\midrule
조직 문화 & 
영업 부서의 AI 결과 불신\cite{davenport2018artificial}; ``왜 거절되었는지 모르겠다''는 내부 반발 \\
\bottomrule
\end{tabularx}
\end{table}


\subsection{거버넌스 체계 설계}
\label{sec:7.1.2}

KBank은 기존 MRM 체계를 확장하여 AI 특화 거버넌스\cite{occ2011sr117}를 구축하였다.

\paragraph{조직 구조}

\begin{figure}[htbp]
\centering
\begin{tikzpicture}[
 node distance=1.5cm,
 box/.style={rectangle, draw, rounded corners, minimum width=3cm, minimum height=0.8cm, align=center, font=\small},
 committee/.style={box, fill=blue!20},
 team/.style={box, fill=green!10},
 unit/.style={box, fill=yellow!10},
]

% 위원회
\node[committee] (board) at (0, 4) {이사회};
\node[committee] (risk) at (0, 2.5) {Risk관리위원회};
\node[committee] (ai) at (0, 1) {AI 거버넌스 위원회};

% 실무 조직
\node[team] (cro) at (-4, 1) {CRO 조직\\(AI 리스크)};
\node[team] (cdo) at (4, 1) {CDO 조직\\(AI 개발)};

% 하위 조직
\node[unit] (mrm) at (-4, -0.5) {ModelRisk관리팀};
\node[unit] (ethics) at (0, -0.5) {AI 윤리팀};
\node[unit] (dev) at (4, -0.5) {AI 개발팀};

% 연결
\draw[->] (board) -- (risk);
\draw[->] (risk) -- (ai);
\draw[->] (ai) -- (cro);
\draw[->] (ai) -- (cdo);
\draw[->] (cro) -- (mrm);
\draw[->] (ai) -- (ethics);
\draw[->] (cdo) -- (dev);

% 점선 연결 (협업)
\draw[dashed, <->] (mrm) -- (ethics);
\draw[dashed, <->] (ethics) -- (dev);

\end{tikzpicture}
\caption{KBank AI 거버넌스 조직 구조}
\label{fig:kbank_org}
\end{figure}

\paragraph{핵심 정책}

KBank이 수립한 AI 신용평가 거버넌스 정책의 핵심 요소:

\begin{enumerate}
 \item \textbf{금지 변수 명확화}: 
 신용정보법 based 금지 변수(성별, 출신 지역, 혼인 여부 등) 목록화 및 
 프록시 의심 변수(우편번호, University명 등) 사용 제한
 
 \item \textbf{공정성 기준 Setup}: 
 4/5 규칙(DP Ratio >= 0.8) 필수 충족; 
 연령대별, 지역별 승인율 분석 의무화
 
 \item \textbf{설명 제공 의무}: 
 모든 거절 건에 대해 상위 4개 영향 요인 제공; 
 SHAP\cite{lundberg2017unified} based 개별 설명 시스템 구축
 
 \item \textbf{인간 검토 기준}: 
 경계선 Case(승인 확률 40-60\%) 및 이의 제기 건 전수 인간 검토
 
 \item \textbf{모니터링\cite{klaise2021monitoring} 주기}: 
 일간 성능 모니터링, 주간 공정성 분석, 분기별 독립 감사
\end{enumerate}


\subsection{공정성 관리 체계}
\label{sec:7.1.3}

KBank의 AI 신용평가 공정성 관리 체계를 상세히 분석한다.

\begin{lstlisting}[language=Python, caption={K-Bank Style Fairness Evaluation Framework}, label={lst:kbank_fairness}]
"""
Financial AI Credit Scoring Fairness Evaluation Framework
Implemented based on K-Bank case (Generalised)
- Prohibited variable check
- Proxy discrimination analysis
- 4/5 rule evaluation
- Adverse Action explanation generation
"""

from dataclasses import dataclass, field
from typing import Dict, List, Optional, Tuple
import numpy as np
import pandas as pd

@dataclass
class CreditFairnessConfig:
 """Credit Scoring Fairness Configuration"""
 
 # Prohibited variables (Based on Credit Information Act)
 prohibited_variables: List[str] = field(default_factory=lambda: [
 'gender', 'sex',
 'birth_region', 'origin_region',
 'marital_status',
 'religion',
 'political_view',
 ])
 
 # Proxy suspect variables
 proxy_suspects: List[str] = field(default_factory=lambda: [
 'zip_code', 'postal_code',
 'university', 'college',
 'employer_industry',
 ])
 
 # Protected attributes for analysis
 protected_attributes: List[str] = field(default_factory=lambda: [
 'age_group', # Young/Middle/Senior
 'region', # Urban/Rural
 'income_level', # Low/Mid/High Income
 ])
 
 # Thresholds
 dp_ratio_threshold: float = 0.8
 proxy_correlation_threshold: float = 0.3


class CreditScoringFairnessEvaluator:
 """
 Credit Scoring AI Fairness Evaluator
 Compliant with FSC guidelines and Credit Information Act
 """
 
 def __init__(self, config: Optional[CreditFairnessConfig] = None):
 self.config = config or CreditFairnessConfig()
 
 def validate_features(
 self,
 feature_names: List[str],
 ) -> Dict[str, List[str]]:
 """Feature validation: Check prohibited and proxy suspect variables"""
 
 prohibited_found = []
 proxy_found = []
 
 for feat in feature_names:
 feat_lower = feat.lower()
 
 # Check prohibited variables
 for prohibited in self.config.prohibited_variables:
 if prohibited in feat_lower:
 prohibited_found.append(feat)
 break
 
 # Check proxy suspect variables
 for proxy in self.config.proxy_suspects:
 if proxy in feat_lower:
 proxy_found.append(feat)
 break
 
 return {
 'prohibited': prohibited_found,
 'proxy_suspects': proxy_found,
 'clean': [f for f in feature_names 
 if f not in prohibited_found and f not in proxy_found],
 }
 
 def analyze_proxy_discrimination(
 self,
 df: pd.DataFrame,
 feature_name: str,
 protected_attribute: str,
 ) -> Dict:
 """Proxy discrimination analysis"""
 
 # Correlation analysis
 if df[feature_name].dtype == 'object':
 # Categorical: Cramer's V
 from scipy.stats import chi2_contingency
 contingency = pd.crosstab(df[feature_name], df[protected_attribute])
 chi2, p_value, dof, expected = chi2_contingency(contingency)
 n = len(df)
 min_dim = min(contingency.shape) - 1
 cramers_v = np.sqrt(chi2 / (n * min_dim)) if min_dim > 0 else 0
 correlation = cramers_v
 else:
 # Numerical: Point Biserial or Spearman
 if df[protected_attribute].nunique() == 2:
 from scipy.stats import pointbiserialr
 # Binary conversion
 binary = pd.factorize(df[protected_attribute])[0]
 corr, p_value = pointbiserialr(binary, df[feature_name])
 correlation = abs(corr)
 else:
 correlation = abs(df[feature_name].corr(
 pd.factorize(df[protected_attribute])[0]
 ))
 
 is_proxy = correlation >= self.config.proxy_correlation_threshold
 
 return {
 'feature': feature_name,
 'protected_attribute': protected_attribute,
 'correlation': correlation,
 'is_proxy': is_proxy,
 'recommendation': 'Remove or Use with Caution' if is_proxy else 'Safe to Use',
 }
 
 def evaluate_approval_rates(
 self,
 y_pred: np.ndarray,
 sensitive_features: pd.DataFrame,
 ) -> Dict[str, Dict]:
 """Approval rate analysis by group (4/5 rule)"""
 
 results = {}
 
 for attr in self.config.protected_attributes:
 if attr not in sensitive_features.columns:
 continue
 
 groups = sensitive_features[attr].unique()
 group_rates = {}
 
 for group in groups:
 mask = sensitive_features[attr] == group
 approval_rate = y_pred[mask].mean()
 group_rates[group] = {
 'approval_rate': approval_rate,
 'count': mask.sum(),
 'approved': y_pred[mask].sum(),
 }
 
 # 4/5 rule evaluation
 rates = [g['approval_rate'] for g in group_rates.values()]
 max_rate = max(rates) if rates else 0
 
 for group, stats in group_rates.items():
 stats['ratio_to_max'] = stats['approval_rate'] / max_rate if max_rate > 0 else 1.0
 stats['passes_4_5'] = stats['ratio_to_max'] >= 0.8
 
 results[attr] = {
 'groups': group_rates,
 'max_rate': max_rate,
 'overall_pass': all(g['passes_4_5'] for g in group_rates.values()),
 'min_ratio': min(g['ratio_to_max'] for g in group_rates.values()),
 }
 
 return results
 
 def generate_adverse_action_notice(
 self,
 shap_values: Dict[str, float],
 feature_descriptions: Dict[str, str],
 customer_name: str = "Customer",
 application_id: str = "",
 ) -> str:
 """
 Generate Adverse Action Notice
 Compliance with Article 36-2 of the Credit Information Act and ECOA Regulation B
 """
 
 # Negative contributing factors (SHAP < 0)
 negative_factors = sorted(
 [(k, v) for k, v in shap_values.items() if v < 0],
 key=lambda x: x[1]
 )[:4] # Top 4
 
 notice = f"""

 Credit Decision Notice


Dear {customer_name},

Application No: {application_id}
Decision Date: {pd.Timestamp.now().strftime('%Y-%m-%d')}

After careful review of your credit application,
we regret to inform you that we are unable to approve it at this time.



 Key Influencing Factors

The following factors primarily influenced this decision:

"""
 
 for i, (feat, _) in enumerate(negative_factors, 1):
 desc = feature_descriptions.get(feat, feat)
 notice += f" {i}. {desc}\n"
 
 notice += """


 Your Rights

1. Request for Further Explanation
 - You may request a detailed explanation of this decision.

2. Right to Appeal
 - If you disagree with this decision, you may appeal within 30 days.

3. Request for Human Review
 - You may request a review by a credit officer.

4. Credit Information Inquiry
 - You may check your credit info once a year for free.



 Contact Us

Customer Center: 1234-5678 (Weekdays 09:00-18:00)
Email: credit@kbank.co.kr
Website: www.kbank.co.kr/credit-decision



This notice is sent in accordance with Article 36-2 of the 
Act on the Use and Protection of Credit Information.

K-Bank Credit Approval Department
"""
 
 return notice
 
 def generate_regulatory_report(
 self,
 model_name: str,
 evaluation_period: str,
 approval_results: Dict,
 total_applications: int,
 ) -> str:
 """Generate report for Financial Supervisory Service submission"""
 
 report = f"""
# AI Credit Scoring Model Fairness Evaluation Report

Model Name: {model_name} 
Evaluation Period: {evaluation_period} 
Total Reviews: {total_applications:,} 
Report Date: {pd.Timestamp.now().strftime('%Y-%m-%d')}

---

## 1. Overview

This report presents the fairness evaluation results for the AI credit scoring model
in compliance with the FSC "AI Guidelines for the Financial Sector" and
the "Act on the Use and Protection of Credit Information".

---

## 2. Approval Rate Analysis by Group

"""
 
 overall_pass = True
 
 for attr, result in approval_results.items():
 attr_pass = result.get('overall_pass', False)
 overall_pass = overall_pass and attr_pass
 
 report += f"### {attr}\n\n"
 report += f"4/5 Rule Met: {'[PASS] Yes' if attr_pass else '[FAIL] No'}\n"
 report += f"Minimum Ratio: {result.get('min_ratio', 0):.4f}\n\n"
 
 report += "| Group | Reviews | Approved | Approval Rate | Ratio to Max | 4/5 Rule |\n"
 report += "|------|----------|----------|--------|--------------|----------|\n"
 
 for group, stats in result.get('groups', {}).items():
 passed = '[PASS]' if stats['passes_4_5'] else '[FAIL]'
 report += f"| {group} | {stats['count']:,} | {stats['approved']:,} | {stats['approval_rate']:.2%} | {stats['ratio_to_max']:.4f} | {passed} |\n"
 
 report += "\n"
 
 report += f"""
---

## 3. Comprehensive Evaluation

| Evaluation Item | Result |
|----------|------|
| Use of Prohibited Variables | [PASS] None |
| Proxy Discrimination Risk | [WARN] Under Monitoring |
| 4/5 Rule Met Overall | {'[PASS] Met' if overall_pass else '[FAIL] Not Met'} |
| Explanation System | [PASS] Operational |

---

## 4. Required Actions

"""
 
 if not overall_pass:
 report += """
### Immediate Action Required

The following actions are applied for groups not meeting the 4/5 rule:

1. Increase human review rate for affected groups (30% -> 50%)
2. Review data augmentation for affected group data during retraining
3. Strengthen monthly fairness monitoring

"""
 else:
 report += """
### Status Quo Maintained

All fairness criteria are currently met.
Maintain regular quarterly monitoring.

"""
 
 report += """
---

## 5. Attachments

1. Detailed Statistical Data (Attachment 1)
2. Model Card (Attachment 2)
3. Fairness Evaluation Methodology (Attachment 3)

---

Author: AI Ethics Team 
Reviewer: Head of Model Risk Management 
Approver: CRO

*Retained for 5 years per Article 15.*
"""
 
 return report
\end{lstlisting}


\subsection{성능 및 교훈}
\label{sec:7.1.4}

\paragraph{주요 성능}

KBank의 AI 신용평가 거버넌스 도입 2년 후 주요 성능:

\begin{table}[htbp]
\centering
\caption{KBank AI 거버넌스 도입 성능}
\label{tab:kbank_outcomes}
\begin{tabularx}{\textwidth}{p{4cm}XX}
\toprule
\textbf{지표} & \textbf{도입 전} & \textbf{도입 2년 후} \\
\midrule
심사 소요 시간 (평균) & 3일 & 10분 (실시간 90\%) \\
연령별 승인율 격차 & 최대 25\%p & 최대 8\%p \\
고객 이의 제기 건수 & 월 150건 & 월 45건 \\
이의 제기 인용률 & 12\% & 8\% \\
금융감독원 지적 사항 & 3건 & 0건 \\
thin-file 고객 승인율 & 15\% & 28\% \\
연체율 (승인 건 기준) & 2.1\% & 2.0\% \\
\bottomrule
\end{tabularx}
\end{table}

\paragraph{핵심 교훈}

\begin{enumerate}
 \item \textbf{기존 MRM 체계 활용}: 
 금융권의 기존 모델 리스크 관리 체계를 AI 거버넌스의 based으로 활용함으로써 
 조직 저항을 최소화하고 빠른 정착이 가능했다.
 
 \item \textbf{설명가능성이 신뢰의 핵심}: 
 영업 부서와 고객 모두 ``왜''에 대한 답변을 원했다. 
 SHAP based 설명 시스템 도입 후 내부 수용도가 급격히 높아졌다.
 
 \item \textbf{프록시 차별 관리의 어려움}: 
 금지 변수는 명확하지만, 프록시 변수는 결정이 어렵다. 
 지속적인 모니터링과 전문가 검토가 필수적이다.
 
 \item \textbf{규제 기관과의 사전 소통}: 
 금융감독원과 AI 모델 도입 전 사전 협의를 통해 
 규제 불확실성을 줄이고 적법성을 확보하였다.
\end{enumerate}

%----------------------------------------------------------
\section{의료 사례: AUniversityHospital의 의료 AI 거버넌스}
\label{sec:7.2}

\subsection{도입 배경}
\label{sec:7.2.1}

AUniversityHospital은 국내 최초로 AI 의료영상 진단 보조\cite{rajpurkar2022ai} 시스템을 
전면 도입한 대형 상급종합Hospital이다. 
흉부 X-ray, CT 폐결절, 유방 촬영술 등 다수의 AI 시스템을 운영하면서 
체계적인 거버넌스 필요성이 대두되었다.

\paragraph{도입 현황 (2025년 기준)}

\begin{table}[htbp]
\centering
\caption{AUniversityHospital AI 의료기기 현황}
\label{tab:hospital_ai_inventory}
\begin{tabularx}{\textwidth}{p{3.5cm}p{2cm}p{2cm}X}
\toprule
\textbf{System} & \textbf{등급} & \textbf{일 처리량} & \textbf{주요 용도} \\
\midrule
흉부 X-ray AI & Class III & 약 2,000건 & 폐렴, 결핵, 폐암 의심 소견 검출 \\
CT 폐결절 AI & Class III & 약 300건 & 폐결절 검출 및 악성도 예측 \\
유방 촬영 AI & Class III & 약 500건 & 유방암 검출 보조 \\
심전도 AI & Class II & 약 1,500건 & 부정맥, 심근경색 조기 경보 \\
안저 영상 AI & Class II & 약 200건 & 당뇨망막병증 선별 \\
\bottomrule
\end{tabularx}
\end{table}


\paragraph{거버넌스 구축 동기}

\begin{itemize}
 \item 다수 AI 시스템의 통합 관리 필요
 \item 의료기기법 및 의료법 준수 체계화
 \item AI 오류로 인한 의료 사고 예방
 \item 의료진의 AI 활용 표준화
 \item 환자 설명 및 동의 절차 정립
\end{itemize}


\subsection{거버넌스 체계}
\label{sec:7.2.2}

AUniversityHospital은 기존 의료기기 관리 위원회를 확장하여 
AI 특화 거버넌스 체계를 구축하였다.

\paragraph{조직 구조}

\begin{itemize}
 \item \textbf{AI 의료기기 관리위원회}: 
 진료부원장(위원장), 의료정보실장, 의료기기관리자, 
 관련 진료과장, 법무팀장, 간호부장으로 구성. 분기별 회의.
 
 \item \textbf{AI 임상검증팀}: 
 각 AI 시스템별 전문의 1인 + 전공의 1인 + 의료정보 연구원 1인으로 구성. 
 월간 성능 검토.
 
 \item \textbf{의료정보실 AI 담당}: 
 AI 시스템 운영, 모니터링, 문서화 담당.
\end{itemize}

\paragraph{핵심 정책}

\begin{enumerate}
 \item \textbf{최종 판독권}: 
 AI 결과는 ``참고 소견''이며, 최종 진단은 반드시 전문의가 수행
 
 \item \textbf{불일치 관리}: 
 AI와 전문의 판독이 불일치하는 경우 전수 기록 및 월간 분석
 
 \item \textbf{환자 고지}: 
 AI 보조 진단 활용 사실을 환자에게 고지 (검사 동의서에 포함)
 
 \item \textbf{하위 집단 모니터링}: 
 연령, 성별, 기저 질환별 성능 분석 수행
 
 \item \textbf{이상사례 보고}: 
 AI 관련 이상사례 발생 시 24시간 내 보고 및 분석
\end{enumerate}


\subsection{임상 검증 및 모니터링}
\label{sec:7.2.3}

\begin{lstlisting}[language=Python, caption={Healthcare AI Clinical Monitoring System}, label={lst:hospital_monitoring}]
"""
Healthcare AI Clinical Monitoring System

Implementation based on A-University Hospital case (Generalized)
- Interpretation Discrepancy Analysis
- Sub-group Performance Monitoring
- Adverse Case Tracking
"""

from dataclasses import dataclass, field
from typing import Dict, List, Optional, Tuple
from datetime import datetime, timedelta
import numpy as np
import pandas as pd
from enum import Enum

class DiscrepancyType(Enum):
 """ days """
 AI_POSITIVE_PHYSICIAN_NEGATIVE = "AI , "
 AI_NEGATIVE_PHYSICIAN_POSITIVE = "AI , "
 SEVERITY_DISAGREEMENT = " Decision days "
 LOCATION_DISAGREEMENT = " Location days "

@dataclass
class ClinicalCase:
 """Clinical """
 case_id: str
 patient_id: str
 modality: str # X-ray, CT, MRI 
 ai_result: Dict # AI physician_result: Dict # final_diagnosis: Optional[Dict] = None # Final (Tracking )
 timestamp: datetime = field(default_factory=datetime.now)
 
 # information ( )
 patient_age: int = 0
 patient_sex: str = ""
 comorbidities: List[str] = field(default_factory=list)


class MedicalAIClinicalMonitor:
 """
 Healthcare AI Clinical Monitoring
 
 Performance Tracking, days Analysis, Monitoring
 """
 
 def __init__(
 self,
 system_name: str,
 sensitivity_threshold: float = 0.90,
 specificity_threshold: float = 0.85,
 ):
 self.system_name = system_name
 self.sensitivity_threshold = sensitivity_threshold
 self.specificity_threshold = specificity_threshold
 
 self.cases: List[ClinicalCase] = []
 self.alerts: List[Dict] = []
 
 def record_case(self, case: ClinicalCase) -> None:
 """ """
 self.cases.append(case)
 
 # days discrepancy = self._check_discrepancy(case)
 if discrepancy:
 self._handle_discrepancy(case, discrepancy)
 
 def _check_discrepancy(self, case: ClinicalCase) -> Optional[DiscrepancyType]:
 """ days """
 
 ai_positive = case.ai_result.get('finding_positive', False)
 physician_positive = case.physician_result.get('finding_positive', False)
 
 if ai_positive and not physician_positive:
 return DiscrepancyType.AI_POSITIVE_PHYSICIAN_NEGATIVE
 elif not ai_positive and physician_positive:
 return DiscrepancyType.AI_NEGATIVE_PHYSICIAN_POSITIVE
 
 return None
 
 def _handle_discrepancy(
 self,
 case: ClinicalCase,
 discrepancy: DiscrepancyType,
 ) -> None:
 """ days """
 
 alert = {
 'case_id': case.case_id,
 'timestamp': datetime.now(),
 'discrepancy_type': discrepancy.value,
 'severity': 'high' if discrepancy == DiscrepancyType.AI_NEGATIVE_PHYSICIAN_POSITIVE else 'medium',
 'ai_result': case.ai_result,
 'physician_result': case.physician_result,
 }
 
 self.alerts.append(alert)
 
 # AI , ( )
 if discrepancy == DiscrepancyType.AI_NEGATIVE_PHYSICIAN_POSITIVE:
 self._escalate_alert(alert)
 
 def _escalate_alert(self, alert: Dict) -> None:
 """ """
 # Implementation: days, SMS, print(f"[CRITICAL] : {alert['case_id']}")
 
 def calculate_performance(
 self,
 start_date: Optional[datetime] = None,
 end_date: Optional[datetime] = None,
 ) -> Dict:
 """Period Performance Calculation"""
 
 # Final evaluated_cases = [
 c for c in self.cases
 if c.final_diagnosis is not None
 ]
 
 if start_date:
 evaluated_cases = [c for c in evaluated_cases if c.timestamp >= start_date]
 if end_date:
 evaluated_cases = [c for c in evaluated_cases if c.timestamp <= end_date]
 
 if not evaluated_cases:
 return {'error': 'No evaluated cases'}
 
 # AI vs Final 
 ai_positive = np.array([c.ai_result.get('finding_positive', False) for c in evaluated_cases])
 final_positive = np.array([c.final_diagnosis.get('confirmed_positive', False) for c in evaluated_cases])
 
 tp = ((ai_positive == True) & (final_positive == True)).sum()
 fn = ((ai_positive == False) & (final_positive == True)).sum()
 fp = ((ai_positive == True) & (final_positive == False)).sum()
 tn = ((ai_positive == False) & (final_positive == False)).sum()
 
 sensitivity = tp / (tp + fn) if (tp + fn) > 0 else 0
 specificity = tn / (tn + fp) if (tn + fp) > 0 else 0
 ppv = tp / (tp + fp) if (tp + fp) > 0 else 0
 npv = tn / (tn + fn) if (tn + fn) > 0 else 0
 
 return {
 'n_cases': len(evaluated_cases),
 'sensitivity': sensitivity,
 'specificity': specificity,
 'ppv': ppv,
 'npv': npv,
 'tp': tp, 'fn': fn, 'fp': fp, 'tn': tn,
 'meets_sensitivity_threshold': sensitivity >= self.sensitivity_threshold,
 'meets_specificity_threshold': specificity >= self.specificity_threshold,
 }
 
 def analyze_subgroups(
 self,
 subgroup_column: str,
 ) -> Dict[str, Dict]:
 """ Performance Analysis"""
 
 evaluated_cases = [c for c in self.cases if c.final_diagnosis is not None]
 
 if subgroup_column == 'age_group':
 # Classification
 def get_age_group(age):
 if age < 40:
 return '40 '
 elif age < 60:
 return '40-59 '
 elif age < 80:
 return '60-79 '
 else:
 return '80 '
 
 groups = {}
 for case in evaluated_cases:
 group = get_age_group(case.patient_age)
 if group not in groups:
 groups[group] = []
 groups[group].append(case)
 
 elif subgroup_column == 'sex':
 groups = {' ': [], ' ': []}
 for case in evaluated_cases:
 if case.patient_sex in ['M', 'male', ' ']:
 groups[' '].append(case)
 elif case.patient_sex in ['F', 'female', ' ']:
 groups[' '].append(case)
 
 else:
 return {'error': f'Unknown subgroup column: {subgroup_column}'}
 
 # Performance Calculation
 results = {}
 for group_name, cases in groups.items():
 if len(cases) < 30: # Minimum 
 results[group_name] = {'n': len(cases), 'insufficient_data': True}
 continue
 
 ai_pos = np.array([c.ai_result.get('finding_positive', False) for c in cases])
 final_pos = np.array([c.final_diagnosis.get('confirmed_positive', False) for c in cases])
 
 tp = ((ai_pos == True) & (final_pos == True)).sum()
 fn = ((ai_pos == False) & (final_pos == True)).sum()
 fp = ((ai_pos == True) & (final_pos == False)).sum()
 tn = ((ai_pos == False) & (final_pos == False)).sum()
 
 sens = tp / (tp + fn) if (tp + fn) > 0 else 0
 spec = tn / (tn + fp) if (tn + fp) > 0 else 0
 
 results[group_name] = {
 'n': len(cases),
 'sensitivity': sens,
 'specificity': spec,
 'positive_rate': final_pos.mean(),
 }
 
 # Analysis
 sensitivities = [r['sensitivity'] for r in results.values() 
 if not r.get('insufficient_data', False)]
 if len(sensitivities) >= 2:
 max_gap = max(sensitivities) - min(sensitivities)
 results['_gap_analysis'] = {
 'max_sensitivity_gap': max_gap,
 'equity_concern': max_gap > 0.10,
 }
 
 return results
 
 def generate_monthly_report(self, month: str) -> str:
 """ Clinical Performance Report Generation"""
 
 perf = self.calculate_performance()
 
 report = f"""
# {self.system_name} Clinical Performance Report

Period: {month} 
 : {perf.get('n_cases', 0):,} 

---

## 1. Performance Metric

| Metric | | | |
|------|-----|------|------|
| (Sensitivity) | {perf.get('sensitivity', 0):.2%} | >={self.sensitivity_threshold:.0%} | {'[PASS]' if perf.get('meets_sensitivity_threshold', False) else '[FAIL]'} |
| (Specificity) | {perf.get('specificity', 0):.2%} | >={self.specificity_threshold:.0%} | {'[PASS]' if perf.get('meets_specificity_threshold', False) else '[FAIL]'} |
| (PPV) | {perf.get('ppv', 0):.2%} | - | - |
| (NPV) | {perf.get('npv', 0):.2%} | - | - |

---

## 2. | | AI | AI |
|--|---------|---------|
| Final | {perf.get('tp', 0)} | {perf.get('fn', 0)} |
| Final | {perf.get('fp', 0)} | {perf.get('tn', 0)} |

---

## 3. Interpretation Discrepancy Analysis

| days | |
|------------|------|
"""
 
 # days discrepancy_counts = {}
 for alert in self.alerts:
 dtype = alert['discrepancy_type']
 discrepancy_counts[dtype] = discrepancy_counts.get(dtype, 0) + 1
 
 for dtype, count in discrepancy_counts.items():
 report += f"| {dtype} | {count} |\n"
 
 report += f"""

---

## 4. Performance

### 
"""
 
 age_analysis = self.analyze_subgroups('age_group')
 report += "| | N | | |\n"
 report += "|--------|---|--------|--------|\n"
 for group, stats in age_analysis.items():
 if group.startswith('_'):
 continue
 if stats.get('insufficient_data', False):
 report += f"| {group} | {stats['n']} | Data | - |\n"
 else:
 report += f"| {group} | {stats['n']} | {stats['sensitivity']:.2%} | {stats['specificity']:.2%} |\n"
 
 report += """

---

## 5. """
 
 if perf.get('fn', 0) > 0:
 report += f"- [WARN] {perf['fn']} Analysis \n"
 
 if not perf.get('meets_sensitivity_threshold', True):
 report += "- [FAIL] : Model \n"
 
 gap_analysis = age_analysis.get('_gap_analysis', {})
 if gap_analysis.get('equity_concern', False):
 report += f"- [WARN] {gap_analysis['max_sensitivity_gap']:.1%}: Equity \n"
 
 report += """

---

 : AI Clinical 
 : 
Approval: 

* Report Healthcare 15 5years .*
"""
 
 return report
\end{lstlisting}


\subsection{성능 및 교훈}
\label{sec:7.2.4}

\paragraph{주요 성능}

\begin{itemize}
 \item 응급 흉부 X-ray AI 도입으로 야간 판독 대기 시간 평균 4시간 -> 30분 단축
 \item 폐결절 AI 활용으로 조기 폐암 발견율 12\% 향상
 \item AI-전문의 불일치 분석을 통한 교육 효과 (전공의 판독 정확도 향상)
 \item 0건의 AI 관련 의료 사고 (거버넌스 도입 후 2년간)
\end{itemize}

\paragraph{핵심 교훈}

\begin{enumerate}
 \item \textbf{전문의 최종 권한의 명확화}: 
 AI는 보조 Tool이며, 최종 결정은 전문의에게 있음을 명확히 함으로써 
 의료진의 수용도를 높이고 법적 책임 소재를 명확히 했다.
 
 \item \textbf{불일치 분석의 가치}: 
 AI와 전문의 판독 불일치는 단순한 오류가 아니라 
 학습 기회로 활용되었다. 이를 통해 AI와 의료진 모두 개선되었다.
 
 \item \textbf{하위 집단 모니터링 필수}: 
 고령 환자에서 AI 성능이 상대적으로 낮음을 발견하고 
 해당 집단에 대한 주의를 강화했다.
 
 \item \textbf{환자 커뮤니케이션}: 
 ``AI가 진단한다''는 오해를 막기 위해 
 ``AI가 의사의 진단을 돕는다''라는 표현을 일관되게 사용했다.
\end{enumerate}

%----------------------------------------------------------
\section{Public 사례: H공공기관의 AI 행정 서비스}
\label{sec:7.3}

\subsection{도입 배경}
\label{sec:7.3.1}

H공공기관은 대규모 인프라를 관리하는 공기업으로, 
AI를 활용한 행정 효율화와 시민 서비스 개선을 추진하였다. 
PublicInstitution 특성상 투명성, 형평성, 민원 Response이 핵심 고려사항이었다.

\paragraph{AI 도입 현황}

\begin{table}[htbp]
\centering
\caption{HPublicInstitution AI 시스템 현황}
\label{tab:public_ai_inventory}
\begin{tabularx}{\textwidth}{p{3cm}p{2cm}X}
\toprule
\textbf{System} & \textbf{리스크 Level} & \textbf{용도} \\
\midrule
민원 자동 분류 & 저위험 & 접수 민원을 담당 부서로 자동 라우팅 \\
\midrule
이용료 감면 심사 Support & 고위험 & 감면 대상 자격 심사 Support (최종 결정은 담당자) \\
\midrule
시설물 이상 탐지 & 중위험 & CCTV, 센서 데이터 based 시설물 이상 감지 \\
\midrule
수요 예측 & 저위험 & 교통량, 이용량 예측을 통한 Resource 배분 \\
\midrule
챗봇 민원 상담 & 저위험 & 단순 문의 자동 응답 \\
\bottomrule
\end{tabularx}
\end{table}


\subsection{거버넌스 체계}
\label{sec:7.3.2}

H공공기관은 행정기본법 제20조와 AI 기본법\cite{cset2025korea_ai_act} 시행령을 준수하기 위한 
체계적인 거버넌스를 구축하였다.

\paragraph{핵심 Principle}

\begin{enumerate}
 \item \textbf{재량 행위 자동화 금지}: 
 행정기본법 제20조에 따라, 담당자 재량이 필요한 결정은 
 AI가 Support하되 최종 결정은 담당자가 수행
 
 \item \textbf{Guarantee of Citizen Rights}: 
 모든 AI based 결정에 대해 설명 요구권, 이의 제기권, 인적 검토권 보장
 
 \item \textbf{알고리즘 투명성}: 
 정보공개법에 따라 알고리즘 정보를 적극적으로 공개
 
 \item \textbf{Equity 모니터링}: 
 지역, 소득 수준, 연령별 서비스 접근성 및 결과 형평성 분석
\end{enumerate}


\paragraph{이용료 감면 심사 AI 거버넌스}

이용료 감면 심사 Support AI는 시민의 경제적 권리에 직접 영향을 미치므로 
가장 엄격한 거버넌스가 적용되었다.

\begin{lstlisting}[language=Python, caption={Public Service AI Governance System}, label={lst:public_governance}]
"""
Public Service AI Governance System

Implementation based on H-Public Institution case (Generalized)
- Discretionary Action Judgment
- Guarantee of Citizen Rights
- Equity Monitoring
"""

from dataclasses import dataclass, field
from typing import Dict, List, Optional, Tuple
from datetime import datetime
from enum import Enum
import pandas as pd

class DecisionType(Enum):
 """ """
 BOUND = " " # DISCRETIONARY = "Discretionary " # Discretionary 

class OversightLevel(Enum):
 """ Level"""
 FULL_AUTO = " " # HUMAN_APPROVAL = " Approval" # Approval AI_SUPPORT = "AI Support" # AI information @dataclass
class CitizenRequest:
 """ """
 request_id: str
 citizen_id: str # 
 request_type: str
 submitted_data: Dict
 timestamp: datetime = field(default_factory=datetime.now)
 
 # information
 ai_recommendation: Optional[Dict] = None
 human_decision: Optional[Dict] = None
 final_decision: Optional[Dict] = None
 
 # information (Equity Analysis , )
 region: str = ""
 age_group: str = ""
 income_level: str = ""

@dataclass
class AppealRequest:
 """ """
 appeal_id: str
 original_request_id: str
 citizen_id: str
 appeal_reason: str
 submitted_at: datetime = field(default_factory=datetime.now)
 
 # review_result: Optional[str] = None
 reviewed_by: Optional[str] = None
 reviewed_at: Optional[datetime] = None


class PublicServiceAIGovernance:
 """
 Public AI Governance
 
 20 , AI Compliance
 """
 
 def __init__(self, service_name: str):
 self.service_name = service_name
 self.requests: List[CitizenRequest] = []
 self.appeals: List[AppealRequest] = []
 
 def classify_decision_type(
 self,
 request_type: str,
 has_discretion: bool,
 ) -> Tuple[DecisionType, OversightLevel]:
 """
 Level Classification
 
 20 Discretionary Automation 
 """
 
 if has_discretion:
 return DecisionType.DISCRETIONARY, OversightLevel.AI_SUPPORT
 else:
 return DecisionType.BOUND, OversightLevel.HUMAN_APPROVAL
 
 def process_request(
 self,
 request: CitizenRequest,
 ai_model,
 decision_type: DecisionType,
 oversight: OversightLevel,
 ) -> Dict:
 """ """
 
 # AI Generation
 ai_recommendation = ai_model.predict(request.submitted_data)
 request.ai_recommendation = {
 'recommendation': ai_recommendation,
 'confidence': ai_model.get_confidence(),
 'explanation': ai_model.get_explanation(),
 'generated_at': datetime.now().isoformat(),
 }
 
 # Level if oversight == OversightLevel.FULL_AUTO:
 # Automation ( )
 if decision_type != DecisionType.BOUND:
 raise ValueError("Discretionary Automation ( 20 )")
 
 request.final_decision = {
 'decision': ai_recommendation,
 'decision_type': 'automated',
 'legal_basis': ' 20 1 ',
 }
 
 else:
 # request.final_decision = {
 'decision': None, # 'decision_type': 'pending_human_review',
 'ai_recommendation': ai_recommendation,
 }
 
 self.requests.append(request)
 
 return {
 'request_id': request.request_id,
 'ai_recommendation': ai_recommendation,
 'oversight_level': oversight.value,
 'requires_human_review': oversight != OversightLevel.FULL_AUTO,
 }
 
 def record_human_decision(
 self,
 request_id: str,
 decision: str,
 decided_by: str,
 override_ai: bool = False,
 override_reason: str = "",
 ) -> None:
 """ """
 
 request = next((r for r in self.requests if r.request_id == request_id), None)
 if not request:
 raise ValueError(f"Request not found: {request_id}")
 
 request.human_decision = {
 'decision': decision,
 'decided_by': decided_by,
 'decided_at': datetime.now().isoformat(),
 'override_ai': override_ai,
 'override_reason': override_reason if override_ai else None,
 }
 
 request.final_decision = {
 'decision': decision,
 'decision_type': 'human_reviewed',
 'ai_recommendation': request.ai_recommendation['recommendation'],
 'ai_followed': not override_ai,
 }
 
 def process_appeal(
 self,
 appeal: AppealRequest,
 ) -> Dict:
 """ """
 
 self.appeals.append(appeal)
 
 # return {
 'appeal_id': appeal.appeal_id,
 'status': 'assigned_for_review',
 'expected_response_days': 14, # 14days 'citizen_rights': [
 ' ',
 'Addition ',
 ' ',
 ],
 }
 
 def analyze_equity(self) -> Dict:
 """Equity Analysis"""
 
 if not self.requests:
 return {'error': 'No data'}
 
 # Analysis
 region_stats = {}
 for request in self.requests:
 if not request.final_decision:
 continue
 
 region = request.region or 'unknown'
 if region not in region_stats:
 region_stats[region] = {'total': 0, 'approved': 0}
 
 region_stats[region]['total'] += 1
 if request.final_decision.get('decision') == 'approved':
 region_stats[region]['approved'] += 1
 
 # Approval Calculation
 for region, stats in region_stats.items():
 stats['approval_rate'] = stats['approved'] / stats['total'] if stats['total'] > 0 else 0
 
 # Analysis
 rates = [s['approval_rate'] for s in region_stats.values() if s['total'] >= 30]
 if len(rates) >= 2:
 max_gap = max(rates) - min(rates)
 equity_concern = max_gap > 0.15 # 15%p else:
 max_gap = 0
 equity_concern = False
 
 return {
 'by_region': region_stats,
 'max_gap': max_gap,
 'equity_concern': equity_concern,
 }
 
 def generate_transparency_report(self) -> str:
 """ Report"""
 
 total_requests = len(self.requests)
 decided_requests = [r for r in self.requests if r.final_decision]
 
 # AI ai_followed = sum(
 1 for r in decided_requests 
 if r.final_decision.get('ai_followed', False)
 )
 ai_follow_rate = ai_followed / len(decided_requests) if decided_requests else 0
 
 # total_appeals = len(self.appeals)
 resolved_appeals = [a for a in self.appeals if a.review_result]
 upheld_appeals = sum(1 for a in resolved_appeals if a.review_result == 'upheld')
 upheld_rate = upheld_appeals / len(resolved_appeals) if resolved_appeals else 0
 
 equity = self.analyze_equity()
 
 report = f"""
# {self.service_name} Report

Period: {datetime.now().strftime('%Yyears %m ')} 
days: {datetime.now().strftime('%Yyears %m %ddays')}

---

## 1. | Items | |
|------|------|
| | {total_requests:,} |
| | {len(decided_requests):,} |
| AI | {ai_follow_rate:.1%} |

---

## 2. | Items | / |
|------|----------|
| | {total_appeals:,} |
| | {len(resolved_appeals):,} |
| | {upheld_rate:.1%} |

---

## 3. Equity Analysis

### Approval 

| | | Approval |
|------|----------|--------|
"""
 
 for region, stats in equity.get('by_region', {}).items():
 if stats['total'] >= 30:
 report += f"| {region} | {stats['total']:,} | {stats['approval_rate']:.1%} |\n"
 
 if equity.get('equity_concern', False):
 report += f"\n[WARN] Approval ({equity['max_gap']:.1%}p) (15%p) . Analysis .\n"
 
 report += f"""

---

## 4. :

1.  : AI based Definition Major .
2.  : 14days .
3.  : .
4. information: Addition information .

---

 : {self.service_name} (1234-5678)

* Report PublicInstitution information .*
"""
 
 return report
\end{lstlisting}


\subsection{성능 및 교훈}
\label{sec:7.3.3}

\paragraph{주요 성능}

\begin{itemize}
 \item 민원 처리 시간 평균 45\% 단축
 \item 민원 분류 정확도 92\% (수동 대비 15\%p 향상)
 \item 이의 제기 인용률 8\% (거버넌스 도입 전 15\%)로 감소 -> 초기 결정 품질 향상 반증
 \item 알고리즘 투명성 Report에 대한 시민 긍정 반응
\end{itemize}

\paragraph{핵심 교훈}

\begin{enumerate}
 \item \textbf{행정기본법 준수의 중요성}: 
 재량행위 자동화 금지 원칙을 명확히 함으로써 법적 Risk를 사전에 차단하였다.
 
 \item \textbf{Guarantee of Citizen Rights이 신뢰의 based}: 
 이의 제기 절차를 쉽게 만들고 투명하게 운영함으로써 
 오히려 이의 제기가 감소하는 역설적 효과가 나타났다.
 
 \item \textbf{Equity Monitoring의 가치}: 
 지역별 승인율 분석을 통해 비수도권 시민에 대한 
 정보 접근성 문제를 발견하고 개선하였다.
 
 \item \textbf{투명성 Report의 소통 효과}: 
 정기적인 알고리즘 투명성 Report 공개가 
 시민 단체 및 언론의 우려를 해소하는 데 효과적이었다.
\end{enumerate}

%----------------------------------------------------------
\section{제조 사례: MHeavy Industries의 AI 품질 관리}
\label{sec:7.4}

\subsection{도입 배경}
\label{sec:7.4.1}

MHeavy Industries은 대형 선박 및 해양 플랜트를 제조하는 기업으로, 
AI 비전 검사 시스템을 도입하여 용접 품질 검사를 Automation하였다. 
제조업 특성상 산업 안전과 품질 보증이 핵심 고려사항이었다.

\paragraph{AI 도입 현황}

\begin{table}[htbp]
\centering
\caption{MHeavy Industries AI 시스템 현황}
\label{tab:manufacturing_ai_inventory}
\begin{tabularx}{\textwidth}{p{3cm}p{2cm}p{2.5cm}X}
\toprule
\textbf{System} & \textbf{리스크 Level} & \textbf{일 처리량} & \textbf{용도} \\
\midrule
용접부 비전 검사 & 고위험 & 약 5,000건 & 용접 결함 자동 검출 \\
\midrule
표면 결함 검사 & 중위험 & 약 3,000건 & 도장 전 표면 결함 검출 \\
\midrule
예측 정비 & 중위험 & 실시간 & 크레인, 절단기 고장 예측 \\
\midrule
안전 모니터링 & 고위험 & 실시간 & 작업자 안전 장구 착용 감지, 위험 구역 침입 탐지 \\
\bottomrule
\end{tabularx}
\end{table}


\subsection{거버넌스 체계}
\label{sec:7.4.2}

MHeavy Industries은 기존 품질 관리 System(QMS)과 안전 관리 System(SMS)을 
AI 거버넌스와 통합하였다.

\paragraph{핵심 Principle}

\begin{enumerate}
 \item \textbf{안전 최우선}: 
 AI 시스템 오류로 인한 안전사고 가능점이 있는 경우 
 보수적으로 결정 (안전 마진 확보)
 
 \item \textbf{품질 보증 연계}: 
 AI 검사 결과는 품질 기록의 일부로 관리되며, 
 고객(선주) 검수 시 제공
 
 \item \textbf{숙련공 검증}: 
 AI 판정에 대한 최종 확인은 숙련 검사원이 수행
 
 \item \textbf{지속적 Learning}: 
 현장 피드백을 반영한 모델 업데이트 체계 구축
\end{enumerate}

\begin{lstlisting}[language=Python, caption={Manufacturing AI Quality Management System}, label={lst:manufacturing_qms}]
"""
Manufacturing AI Quality Management System

Implementation based on M-Heavy Industries case (Generalized)
- Welding Defect Inspection
- Quality Record Management
- Skilled Worker Verification System
"""

from dataclasses import dataclass, field
from typing import Dict, List, Optional, Tuple
from datetime import datetime
from enum import Enum
import numpy as np

class DefectSeverity(Enum):
 """ ( )"""
 NONE = " None"
 MINOR = " ( )"
 MAJOR = " ( )"
 CRITICAL = " ( )"

class InspectionResult(Enum):
 """ """
 PASS = " "
 CONDITIONAL_PASS = " "
 FAIL = " "
 REINSPECTION = " "

@dataclass
class WeldInspection:
 """ """
 inspection_id: str
 work_order: str # weld_location: str # Location Code
 welder_id: str # ID
 
 # AI ai_result: Dict = field(default_factory=dict)
 ai_confidence: float = 0.0
 ai_defects_detected: List[Dict] = field(default_factory=list)
 
 # Skilled Worker Verification
 inspector_id: Optional[str] = None
 inspector_result: Optional[Dict] = None
 inspector_override: bool = False
 override_reason: str = ""
 
 # Final 
 final_result: Optional[InspectionResult] = None
 
 timestamp: datetime = field(default_factory=datetime.now)


class ManufacturingQualityAI:
 """
 Manufacturing AI Management
 
 Welding Defect Inspection, Quality Record Management
 """
 
 def __init__(
 self,
 system_name: str,
 min_detection_rate: float = 0.99, # Minimum 
 max_overkill_rate: float = 0.05, # ):
 self.system_name = system_name
 self.min_detection_rate = min_detection_rate
 self.max_overkill_rate = max_overkill_rate
 
 self.inspections: List[WeldInspection] = []
 
 def classify_defect_severity(
 self,
 defect_type: str,
 defect_size_mm: float,
 weld_class: str, # A, B, C ( )
 ) -> DefectSeverity:
 """ Classification ( )"""
 
 # (Example)
 thresholds = {
 'A': {'minor': 2.0, 'major': 4.0}, # Class A: 'B': {'minor': 3.0, 'major': 6.0},
 'C': {'minor': 5.0, 'major': 10.0},
 }
 
 thresh = thresholds.get(weld_class, thresholds['B'])
 
 if defect_size_mm <= thresh['minor']:
 return DefectSeverity.MINOR
 elif defect_size_mm <= thresh['major']:
 return DefectSeverity.MAJOR
 else:
 return DefectSeverity.CRITICAL
 
 def process_inspection(
 self,
 inspection: WeldInspection,
 ) -> Dict:
 """ """
 
 # AI based 
 if inspection.ai_confidence >= 0.95:
 # : AI based Classification
 if not inspection.ai_defects_detected:
 routing = "auto_pass" # else:
 max_severity = max(
 (self.classify_defect_severity(
 d['type'], d['size_mm'], d.get('weld_class', 'B')
 ) for d in inspection.ai_defects_detected),
 default=DefectSeverity.NONE
 )
 if max_severity == DefectSeverity.CRITICAL:
 routing = "immediate_review" # else:
 routing = "standard_review" # days else:
 # : routing = "low_confidence_review"
 
 self.inspections.append(inspection)
 
 return {
 'inspection_id': inspection.inspection_id,
 'ai_confidence': inspection.ai_confidence,
 'defects_count': len(inspection.ai_defects_detected),
 'routing': routing,
 'requires_human_review': routing != "auto_pass",
 }
 
 def record_inspector_decision(
 self,
 inspection_id: str,
 inspector_id: str,
 result: InspectionResult,
 override_ai: bool = False,
 override_reason: str = "",
 additional_defects: List[Dict] = None,
 ) -> None:
 """Skilled Worker Verification """
 
 inspection = next(
 (i for i in self.inspections if i.inspection_id == inspection_id),
 None
 )
 if not inspection:
 raise ValueError(f"Inspection not found: {inspection_id}")
 
 inspection.inspector_id = inspector_id
 inspection.inspector_result = {
 'result': result.value,
 'override_ai': override_ai,
 'override_reason': override_reason,
 'additional_defects': additional_defects or [],
 'timestamp': datetime.now().isoformat(),
 }
 inspection.inspector_override = override_ai
 inspection.override_reason = override_reason
 inspection.final_result = result
 
 def calculate_performance(self) -> Dict:
 """AI Performance Calculation"""
 
 # Skilled Worker Verification verified = [i for i in self.inspections if i.inspector_result]
 
 if not verified:
 return {'error': 'No verified inspections'}
 
 # AI vs Comparison
 ai_detected_defect = np.array([
 len(i.ai_defects_detected) > 0 for i in verified
 ])
 inspector_confirmed_defect = np.array([
 i.final_result in [InspectionResult.FAIL, InspectionResult.CONDITIONAL_PASS]
 for i in verified
 ])
 
 tp = ((ai_detected_defect == True) & (inspector_confirmed_defect == True)).sum()
 fn = ((ai_detected_defect == False) & (inspector_confirmed_defect == True)).sum() # Escape
 fp = ((ai_detected_defect == True) & (inspector_confirmed_defect == False)).sum() # Overkill
 tn = ((ai_detected_defect == False) & (inspector_confirmed_defect == False)).sum()
 
 detection_rate = tp / (tp + fn) if (tp + fn) > 0 else 1.0
 overkill_rate = fp / (fp + tn) if (fp + tn) > 0 else 0.0
 escape_rate = fn / (tp + fn) if (tp + fn) > 0 else 0.0
 
 # Override Analysis
 override_count = sum(1 for i in verified if i.inspector_override)
 override_rate = override_count / len(verified) if verified else 0
 
 return {
 'n_inspections': len(verified),
 'detection_rate': detection_rate,
 'escape_rate': escape_rate,
 'overkill_rate': overkill_rate,
 'override_rate': override_rate,
 'meets_detection_threshold': detection_rate >= self.min_detection_rate,
 'meets_overkill_threshold': overkill_rate <= self.max_overkill_rate,
 'confusion_matrix': {'tp': tp, 'fn': fn, 'fp': fp, 'tn': tn},
 }
 
 def generate_quality_report(self, period: str) -> str:
 """ Report Generation"""
 
 perf = self.calculate_performance()
 
 report = f"""
# {self.system_name} AI Performance Report

Period: {period} 
 : {perf.get('n_inspections', 0):,} 

---

## 1. Performance Metric

| Metric | | | |
|------|-----|------|------|
| | {perf.get('detection_rate', 0):.2%} | >={self.min_detection_rate:.0%} | {'[PASS]' if perf.get('meets_detection_threshold', False) else '[FAIL]'} |
| Escape Rate | {perf.get('escape_rate', 0):.2%} | Minimum | - |
| Overkill Rate | {perf.get('overkill_rate', 0):.2%} | <={self.max_overkill_rate:.0%} | {'[PASS]' if perf.get('meets_overkill_threshold', False) else '[WARN]'} |
| Override | {perf.get('override_rate', 0):.2%} | | - |

---

## 2. | | AI | AI |
|--|-------------|---------------|
|   | {perf.get('confusion_matrix', {}).get('tp', 0)} | {perf.get('confusion_matrix', {}).get('fn', 0)} (Escape) |
|   | {perf.get('confusion_matrix', {}).get('fp', 0)} (Overkill) | {perf.get('confusion_matrix', {}).get('tn', 0)} |

---

## 3. Impact Analysis

"""
 
 escape_count = perf.get('confusion_matrix', {}).get('fn', 0)
 if escape_count > 0:
 report += f"""
[WARN] Escape: {escape_count} 

AI .
 Analysis .

"""
 else:
 report += "[PASS] Escape 0 : AI 1 .\n\n"
 
 report += f"""
---

## 4. """
 
 if not perf.get('meets_detection_threshold', True):
 report += "- [FAIL] : Model Learning \n"
 
 if perf.get('overkill_rate', 0) > self.max_overkill_rate:
 report += "- [WARN] Overkill : ( , )\n"
 
 if perf.get('override_rate', 0) > 0.15:
 report += f"- [WARN] Override {perf['override_rate']:.1%}: AI- days Analysis \n"
 
 report += """

---

 : Management 
 : QA 
Approval: 

* Report .*
* 15 5years .*
"""
 
 return report
\end{lstlisting}


\subsection{성능 및 교훈}
\label{sec:7.4.3}

\paragraph{주요 성능}

\begin{itemize}
 \item 용접 검사 시간 60\% 단축 (검사원 1인당 처리량 2.5배 증가)
 \item Escape Rate 0.1\% 이하 유지 (업계 평균 0.5\% 대비 우수)
 \item 재작업률 15\% 감소 (조기 결함 발견으로 인한 비용 절감)
 \item AI 관련 품질 클레임 0건 (도입 후 3년간)
\end{itemize}

\paragraph{핵심 교훈}

\begin{enumerate}
 \item \textbf{안전 마진 확보}: 
 검출율을 최우선으로 하고, Overkill은 어느 정도 감수하는 전략이 
 품질 사고 예방에 효과적이었다.
 
 \item \textbf{숙련공 협업 필수}: 
 AI를 숙련공 대체가 아닌 Support Tool로 포지셔닝함으로써 
 현장 수용도를 높이고 지속적인 피드백을 확보하였다.
 
 \item \textbf{품질 기록 통합된}: 
 AI 검사 결과를 기존 품질 관리 시스템에 통합하여 
 추적성(Traceability)을 확보하고 선급 검사 Response이 용이해졌다.
 
 \item \textbf{지속적 학습 체계}: 
 숙련공 Override 사례를 학습 데이터로 활용하여 
 모델 성능을 지속적으로 개선하는 피드백 루프를 구축하였다.
\end{enumerate}


%----------------------------------------------------------
\subsection*{제7장 요약}

본 장에서는 한국의 금융, 의료, 공공, 제조 분야에서의 
AI 거버넌스 구축 사례를 분석하였다. 주요 시사점을 요약하면 다음과 같다.

\begin{enumerate}
 \item \textbf{기존 체계 활용}: 
 모든 사례에서 기존 리스크 관리, 품질 관리, 규제 준수 체계를 
 AI 거버넌스의 based으로 활용하여 효과적인 정착이 가능했다.
 
 \item \textbf{인간-AI 협업 설계}: 
 AI를 대체가 아닌 Support Tool로 포지셔닝하고, 
 최종 결정권은 인간에게 유지함으로써 
 내부 수용도와 법적 안정성을 확보하였다.
 
 \item \textbf{설명과 투명성}: 
 모든 분야에서 ``왜''에 대한 답변이 핵심이었다. 
 고객, 환자, 시민, 숙련공 모두 AI 결정 근거를 알고 싶어 했다.
 
 \item \textbf{지속적 모니터링}: 
 배포 후 지속적인 성능, 공정성, 안전성 모니터링(\ref{sec:4.3.4})이 
 문제의 조기 발견과 신뢰 유지에 필수적이었다.
 
 \item \textbf{산업 특수성 반영}: 
 보편적인 거버넌스 원칙을 산업 특성에 맞게 적용하는 것이 중요했다. 
 금융의 공정성, 의료의 안전성, 공공의 형평성, 제조의 품질이 
 각각 핵심 가치였다.
\end{enumerate}

이러한 사례들은 이론적 프레임워크가 실제 현장에서 
어떻게 적용될 수 있는지를 보여주며, 
유사한 상황에 있는 조직들에게 실질적인 참고가 될 것이다. 
마지막 제8장에서는 AI 기술의 비약적인 발전과 규제 환경의 변화 속에서 
AI 거버넌스가 나아가야 할 미래 방향을 제시하며 본 책을 마무리한다.
